% =============================================================================
% 5.1 Requisitos Funcionales del Sistema
% Archivo: 5.1_requisitos_funcionales.tex
% =============================================================================

\section{Requisitos Funcionales del Sistema}
\label{sec:requisitos-funcionales}

La definici\'on formal de los requisitos funcionales constituye una etapa cr\'itica
en el ciclo de vida del desarrollo de software, puesto que establece de manera
inequ\'ivoca las capacidades que el sistema debe proveer para satisfacer las
necesidades operativas de las organizaciones no gubernamentales. El presente
cat\'alogo recopila la totalidad de los requisitos funcionales identificados
durante las fases de descubrimiento, an\'alisis y validaci\'on del \textit{Design
Documentation System} (DDS), los cuales fueron refinados iterativamente con
base en el an\'alisis documental del c\'odigo fuente existente (ingenier\'ia
inversa) y el estudio de los procesos organizacionales reales.

Cada requisito funcional se encuentra codificado mediante un identificador
\'unico con el prefijo \texttt{RF-}, acompa\~nado de un nombre descriptivo y una
especificaci\'on concisa del comportamiento esperado. Esta estructura facilita
la trazabilidad bidireccional entre los requisitos, los casos de prueba y los
m\'odulos implementados en la arquitectura del sistema. La Tabla~\ref{tab:requisitos-funcionales}
presenta el cat\'alogo completo de veintinueve (29) requisitos funcionales que
rigen el dise\~no y la construcci\'on de la plataforma.

\begin{small}
\begin{longtable}{|p{1.4cm}|p{3.8cm}|p{8.5cm}|}
\caption{Cat\'alogo de requisitos funcionales del sistema}
\label{tab:requisitos-funcionales} \\
\hline
\textbf{C\'odigo} & \textbf{Nombre} & \textbf{Descripci\'on} \\
\hline
\endfirsthead

\multicolumn{3}{c}%
{{\tablename\ \thetable{} -- Continuaci\'on de la p\'agina anterior}} \\
\hline
\textbf{C\'odigo} & \textbf{Nombre} & \textbf{Descripci\'on} \\
\hline
\endhead

\hline
\multicolumn{3}{r}{{Contin\'ua en la siguiente p\'agina\ldots}} \\
\endfoot

\hline
\endlastfoot

RF-001 & Validar RUC de organizaci\'on & El sistema debe validar el RUC contra SUNAT antes del registro. \\
\hline
RF-002 & Registrar organizaci\'on (Tenant) & El sistema debe crear el entorno aislado y la cuenta administradora. \\
\hline
RF-003 & Evaluar riesgo de acceso & El sistema debe evaluar el nivel de riesgo en cada inicio de sesi\'on. \\
\hline
RF-004 & Verificar c\'odigo OTP & El sistema debe validar los c\'odigos de seguridad enviados por correo. \\
\hline
RF-005 & Reenviar c\'odigo OTP & El sistema debe permitir reenviar los c\'odigos de seguridad si no llegaron. \\
\hline
RF-006 & Autenticar por terminal & El sistema debe permitir el acceso r\'apido mediante un c\'odigo PIN num\'erico. \\
\hline
RF-007 & Gestionar roles del sistema & El sistema debe permitir crear, editar y eliminar roles personalizados. \\
\hline
RF-008 & Asignar permisos a roles & El sistema debe permitir asignar capacidades y permisos espec\'ificos a cada rol. \\
\hline
RF-009 & Asignar roles a usuarios & El sistema debe permitir vincular usuarios con roles en sedes f\'isicas espec\'ificas. \\
\hline
RF-010 & Supervisar sesiones activas & El sistema debe permitir visualizar y revocar conexiones activas de los usuarios. \\
\hline
RF-011 & Gestionar sedes & El sistema debe permitir registrar, editar y desactivar sedes o locales. \\
\hline
RF-012 & Gestionar voluntarios & El sistema debe permitir registrar y administrar perfiles completos de voluntarios. \\
\hline
RF-013 & Gestionar beneficiarios & El sistema debe permitir registrar beneficiarios utilizando perfiles m\'edicos y familiares diferenciados. \\
\hline
RF-014 & Auditar acceso m\'edico & El sistema debe registrar obligatoriamente el motivo al acceder a datos m\'edicos. \\
\hline
RF-015 & Emitir carnets digitales & El sistema debe permitir dise\~nar y generar carnets de identificaci\'on con c\'odigos QR. \\
\hline
RF-016 & Gestionar admisi\'on & El sistema debe administrar el ciclo de selecci\'on y entrevistas de voluntarios. \\
\hline
RF-017 & Autoregistrar candidatos & El sistema debe generar enlaces p\'ublicos para que los candidatos se inscriban solos. \\
\hline
RF-018 & Gestionar proyectos & El sistema debe permitir planificar proyectos, definir fechas y asignar presupuestos. \\
\hline
RF-019 & Gestionar tareas y actividades & El sistema debe permitir dividir proyectos en tareas y programar actividades presenciales. \\
\hline
RF-020 & Asignar voluntarios a actividades & El sistema debe permitir designar qu\'e voluntarios participar\'an en cada actividad programada. \\
\hline
RF-021 & Registrar asistencia y horas & El sistema debe permitir marcar la participaci\'on y recolectar evidencias por actividad. \\
\hline
RF-022 & Gestionar inventario de art\'iculos & El sistema debe permitir administrar el cat\'alogo de art\'iculos y materiales disponibles. \\
\hline
RF-023 & Registrar movimientos de inventario & El sistema debe registrar las entradas, salidas y transferencias de cada art\'iculo. \\
\hline
RF-024 & Gestionar cuentas financieras & El sistema debe permitir registrar cuentas bancarias y contabilizar todos los ingresos. \\
\hline
RF-025 & Aprobar egresos financieros & El sistema debe exigir un flujo de aprobaci\'on interno para autorizar gastos. \\
\hline
RF-026 & Gestionar notificaciones & El sistema debe permitir configurar plantillas de correo y revisar mensajes enviados. \\
\hline
RF-027 & Auditar historial del sistema & El sistema debe registrar inalterablemente cualquier modificaci\'on realizada por los usuarios. \\
\hline
RF-028 & Restringir accesos por contexto & El sistema debe permitir bloquear el acceso seg\'un el horario o direcci\'on IP. \\
\hline
RF-029 & Resumir m\'etricas de seguridad & El sistema debe usar IA para explicar reportes de seguridad y actividades sospechosas. \\

\end{longtable}
\end{small}

Los requisitos funcionales presentados abarcan siete dominios operativos
fundamentales: (a)~autenticaci\'on y seguridad adaptativa (RF-001 a RF-006),
(b)~gesti\'on de identidad y acceso (RF-007 a RF-010), (c)~administraci\'on
organizacional (RF-011 a RF-015), (d)~gesti\'on del voluntariado y admisi\'on
(RF-016 a RF-017), (e)~planificaci\'on y ejecuci\'on de proyectos (RF-018 a RF-021),
(f)~gesti\'on de recursos e inventario (RF-022 a RF-025), y (g)~auditor\'ia,
notificaciones e inteligencia de seguridad (RF-026 a RF-029). Esta organizaci\'on
modular refleja la arquitectura multi-inquilino del sistema y garantiza que cada
dominio funcional pueda evolucionar de forma independiente sin comprometer la
integridad del conjunto.

Cabe destacar que los requisitos RF-003, RF-014, RF-027 y RF-029 incorporan
mecanismos avanzados de seguridad y auditor\'ia que trascienden las
funcionalidades convencionales de los sistemas de gesti\'on de ONG, lo cual
constituye uno de los aportes diferenciadores de la presente investigaci\'on.
En particular, el requisito RF-029 integra capacidades de inteligencia artificial
generativa para la s\'intesis de m\'etricas de seguridad, alineando la plataforma
con las tendencias emergentes en ciberseguridad adaptativa.
